\subsection*{\S{} 5. 域的伽罗瓦扩张环。结构定理。}

\paragraph{1. 基础环。}
设 $K$ 是一个\emph{基域 $P$ 上的 $n$ 次域}，并且假定 $K$ 是可分扩张。设 $\Gamma$ 是相应的伽罗瓦域；它在同构意义下唯一确定，包含 $n$ 个共轭域 $K^{(i)}$，并且是它们的直积，其中规定 $K=K^{(1)}$。用 $\mathfrak K$ 表示 $P$ 的一个扩张，它同 $K$、因而也同所有 $K^{(i)}$ 同构，并且相对于 $P$ 与它们等价；此外还要求其与 $\Gamma$ 的交 $[\mathfrak K,\Gamma]$ 就是 $P$（亦即 $\Gamma$ 与 $\mathfrak K$ 不存在共同的包络域）。\srcfn{15)}{这里当然假定——这并不限制一般性——$K$ 与 $\mathfrak K$ 中不属于 $P$ 的元素之间不存在任何关系（参见第 1 节第 2 小节、第 2 节第 1 小节以及脚注 11）。}由于 $\mathfrak K$ 具有一个独立、有限且包含单位元的 $P$-模基，根据第 2 节第 1 小节，直积 $\mathfrak K_\Gamma$、$\mathfrak K_K$、$\mathfrak K_{K^{(i)}}$ 都存在；$\mathfrak K_\Gamma$ 是它的 $n$ 个子环 $\mathfrak K_{K^{(i)}}$ 的直积。称 $\mathfrak K_\Gamma$ 为 $\mathfrak K$ 的\emph{伽罗瓦扩张环}。同时，$\mathfrak K_\Gamma$ 还是 $\mathfrak K_K$ 与 $\Gamma$ 关于 $K$ 的直积，因为 $\mathfrak K$ 的每个独立 $P$-基也成为 $\mathfrak K_K$ 的独立 $K$-基；对各共轭域也同样如此。\srcfn{16)}{对所有共轭域一致成立的定理，通常只对 $\mathfrak K_K$ 表述，也就是省去指标。}这些环的理想之间有如下关系：

\paragraph{定理。}
\emph{$\mathfrak K$ 中的每个理想，都是它在 $\mathfrak K_K$ 或 $\mathfrak K_\Gamma$ 中的扩张的收缩理想；$\mathfrak K_K$ 中的每个理想，都是它在 $\mathfrak K_\Gamma$ 中的扩张的收缩理想。}

\paragraph{证明。}
由于 $P$ 和 $K$ 都是域，$\mathfrak K$ 中的每个 $P$-模以及 $\mathfrak K_K$ 中的每个 $K$-模，都分别具有一个独立有限的 $P$-模基或 $K$-模基，而且该基可以分别补充成 $\mathfrak K$ 或 $\mathfrak K_K$ 的这种基。于是该定理由第 2 节第 2 小节中的定理推出。

用 $\mathfrak B_{K^{(i)}}^{(i)}$ 表示 $\mathfrak K_{K^{(i)}}$ 的差分理想，即
\[
\mathfrak B_{K^{(i)}}^{(i)}=\{\ldots,x-\xi^{(i)},\ldots\},
\]
特别地：
\[
\mathfrak B_K=\{\ldots,x-\xi,\ldots\}.
\]
由于 $\xi$ 或 $\xi^{(i)}$ 分别是 $K$ 或 $K^{(i)}$ 中通过同构与 $x$ 对应的元素，所以这 $n$ 个差分理想彼此共轭。

用 $\mathfrak A_K$ 或 $\mathfrak A_{K^{(i)}}^{(i)}$ 表示 $\mathfrak K_K$ 中的差分商 $(0):\mathfrak B_K$ 及其各共轭。它们在 $\mathfrak K_\Gamma$ 中的扩张理想分别记作 $\mathfrak B_\Gamma$ 和 $\mathfrak A_\Gamma$。由本小节关于扩张理想与收缩理想的定理，特别得到：

\emph{差分理想 $\mathfrak B_K$ 和差分商 $\mathfrak A_K$，分别是它们的扩张 $\mathfrak B_\Gamma$ 和 $\mathfrak A_\Gamma$ 的收缩理想。}

\paragraph{2. 伽罗瓦扩张环的直和分解。}
差分理想及差分商同伽罗瓦扩张环结构之间的关系，由下述定理给出。

\paragraph{结构定理。}
\emph{伽罗瓦扩张环 $\mathfrak K_\Gamma$ 的零理想，是 $n$ 个共轭差分理想的扩张之交；$\mathfrak K_\Gamma$ 本身，是 $n$ 个共轭差分商的扩张之直和。$\mathfrak K_{K^{(i)}}$ 是相应差分理想与差分商的直和，而它的零理想是二者之交。}

\paragraph{证明。}
按每个 $\mathfrak B_\Gamma^{(i)}$ 取的剩余类环关于 $\Gamma$ 的秩为一，因为每个元素模 $\mathfrak B_\Gamma^{(i)}$ 都同 $\Gamma$ 中的一个元素同余，而 $\mathfrak B_\Gamma^{(i)}$ 不可能包含 $\Gamma$ 中任何非零元素（因为令 $x\to\xi^{(i)}$ 时 $\mathfrak B_\Gamma^{(i)}$ 消失）。由于可分性，各 $\mathfrak B_\Gamma^{(i)}$ 全都不同，因而它们两两互素（任意两个之和等于 $\mathfrak K_\Gamma$）。因此，理想
\[
\mathfrak B_\Gamma^{(1)},
[\mathfrak B_\Gamma^{(1)},\mathfrak B_\Gamma^{(2)}],\ldots,
[\mathfrak B_\Gamma^{(1)},\ldots,\mathfrak B_\Gamma^{(n)}]
\]
的秩每一步至少降低一个单位；所有这些理想之交成为零理想，而由于两两互素，少于全体共轭理想的交必定异于零理想（所以秩在每一步恰好降低一个单位）。由熟知的定理，这便给出 $\mathfrak K_\Gamma$ 如下唯一确定的直和表示：\srcfn{17)}{例如参见 E. Noether，\emph{Abstrakter Aufbau der Idealtheorie in algebraischen Zahl- und Funktionenkörpern}，Math. Ann. 96 (1927)，26--61，第 4 节第 5 小节。}
\begin{align}
\mathfrak K_\Gamma
&=\mathfrak C_\Gamma^{(1)}+\cdots+\mathfrak C_\Gamma^{(n)}
=\Gamma e^{(1)}+\cdots+\Gamma e^{(n)}
=\mathfrak C_\Gamma^{(i)}+\mathfrak B_\Gamma^{(i)}, \tag{1}\\
e&=e^{(1)}+\cdots+e^{(n)},\qquad
\mathfrak C_\Gamma^{(i)}=
[\mathfrak B_\Gamma^{(1)},\ldots,\mathfrak B_\Gamma^{(i-1)},
\mathfrak B_\Gamma^{(i+1)},\ldots,\mathfrak B_\Gamma^{(n)}].\notag
\end{align}
还须证明 $\mathfrak C_\Gamma$ 等于 $\mathfrak A_\Gamma$。只要证明 $\mathfrak K_K$ 中相应的二者相等，这一结论便随之成立。在这里有分解：
\begin{equation}
\mathfrak K_K=\mathfrak C_K+\mathfrak B_K
=\mathfrak K_K e^{(1)}+\mathfrak K_K(e-e^{(1)}). \tag{2}
\end{equation}
事实上，$e^{(1)}$ 属于 $\mathfrak K_K$；因为 $\mathfrak C_\Gamma^{(1)}$ 在 $\mathfrak B_\Gamma^{(2)},\ldots,\mathfrak B_\Gamma^{(n)}$ 的所有置换下保持不变，所以作为 $\mathfrak C_\Gamma^{(1)}$ 中单位元分量的 $e^{(1)}$ 也同样保持不变。因此，用 $\mathfrak K_K$ 的一个独立 $K$-基来表示 $e^{(1)}$——这是可能的，因为该基同时也是 $\mathfrak K_\Gamma$ 的独立 $\Gamma$-基——则各个系数在 $K^{(2)},\ldots,K^{(n)}$ 的所有置换下都保持不变，因而属于 $K$。由于 $(e^{(1)})^2=e^{(1)}$，和式表示
\[
\mathfrak K_K=\mathfrak K_K e^{(1)}+\mathfrak K_K(e-e^{(1)})
\]
是直和；它在 $\mathfrak K_\Gamma$ 中的扩张仍然是直和，故由 (1) 得到 $\mathfrak K_\Gamma=\mathfrak C_\Gamma^{(1)}+\mathfrak B_\Gamma^{(1)}$。于是根据第 2 节第 2 小节的末尾部分，有
\[
\mathfrak K_K e^{(1)}=[\mathfrak C_\Gamma,\mathfrak K_K]
=[\Gamma e^{(1)},\mathfrak K_K]=Ke^{(1)}=\mathfrak C_K;
\]
\[
\mathfrak K_K(e-e^{(1)})=[\mathfrak B_\Gamma,\mathfrak K_K]=\mathfrak B_K,
\]
从而 (2) 得证。

由 (2) 又可推出 $\mathfrak C_K=\mathfrak A_K$，进而 $\mathfrak C_\Gamma=\mathfrak A_\Gamma$。事实上，由 $e^{(1)}(e-e^{(1)})=0$ 可得 $\mathfrak C_K\mathfrak B_K=0$；反过来，$b\mathfrak B_K=0$ 也给出 $b(e-e^{(1)})=0$，所以 $b=be=be^{(1)}$，于是
\[
\mathfrak C_K=(0):\mathfrak B_K=\mathfrak A_K,
\]
并且扩张后有 $\mathfrak C_\Gamma=\mathfrak A_\Gamma$。因此：
\begin{align}
\mathfrak K_\Gamma
&=\mathfrak A_\Gamma^{(1)}+\cdots+\mathfrak A_\Gamma^{(n)}
=\Gamma e^{(1)}+\cdots+\Gamma e^{(n)};
&0&=[\mathfrak B_\Gamma^{(1)},\ldots,\mathfrak B_\Gamma^{(n)}]; \tag{3}\\
\mathfrak K_{K^{(i)}}
&=\mathfrak A_{K^{(i)}}^{(i)}+\mathfrak B_{K^{(i)}}^{(i)}
=K^{(i)}e^{(i)}+\mathfrak B_{K^{(i)}}^{(i)};
&0&=[\mathfrak A_{K^{(i)}}^{(i)},\mathfrak B_{K^{(i)}}^{(i)}].\notag
\end{align}
至此，该定理的各部分全都得证。还应指出，这是一个\emph{完全不变的结构定理}，因为其中出现的所有构造（$\mathfrak K_K$、$\mathfrak K_\Gamma$、$\mathfrak B_K$、$\mathfrak A_K$）都不依赖于基来定义。\srcfn{18)}{这里也可以借助定义方程以及相应多项式的线性因式分解来表述，但这样的表述并非不变的。}

\paragraph{3. 由共轭元素给出的 $\mathfrak K$ 的分量表示。}
与共轭差分理想及差分商相对应，$\mathfrak K$ 的元素可通过共轭分量来表示：

\paragraph{定理。}
\emph{设
\[
x=\xi^{(1)}e^{(1)}+\cdots+\xi^{(n)}e^{(n)}
\]
是 $\mathfrak K$ 中元素 $x$ 的 $\mathfrak K$-分量表示，则这 $n$ 个分量彼此共轭。特别地，$\xi^{(1)},\ldots,\xi^{(n)}$ 正是通过 $\mathfrak K$ 到 $K^{(1)},\ldots,K^{(n)}$ 的各同构与 $x$ 对应的 $n$ 个共轭值。}

$e^{(1)},\ldots,e^{(n)}$ 彼此共轭，可由第 2 小节中的关系式 (3) 看出：把 $\mathfrak B_K$ 映到 $\mathfrak B_{K^{(i)}}^{(i)}$ 的 $\Gamma$ 的自同构，也同时把 $e^{(1)}$ 映到 $e^{(i)}$，而后者确实是 $\mathfrak K_{K^{(i)}}$ 中的元素。同一组关系式还给出 $xe^{(i)}=\xi^{(i)}e^{(i)}$；这是因为 $\mathfrak B_{K^{(i)}}^{(i)}e^{(i)}=0$，从而
\[
(x-\xi^{(i)})e^{(i)}=0
\]
（也可由以下事实推出：作为域的 $\mathfrak K$ 同它的分量同构，而该分量就是 $K^{(i)}e^{(i)}$，其中 $x$ 与 $\xi^{(i)}e^{(i)}$ 同构对应）。

由该定理可得：

\paragraph{推论。}
\emph{设 $\mathfrak B$ 是 $\mathfrak K$ 中的一个绝对模（即它在包含任意两个元素的同时也包含二者之差），并设 $\mathfrak B^{(1)},\ldots,\mathfrak B^{(n)}$ 是它在 $\mathfrak K_\Gamma$ 中的各分量，则 $\mathfrak B$ 等于交 $[\mathfrak K,\mathfrak B^{(1)}+\cdots+\mathfrak B^{(n)}]$。}

事实上，$\mathfrak B$ 包含在这个交中；反过来，若 $\mathfrak B^{(1)}+\cdots+\mathfrak B^{(n)}$ 中的一个元素还属于 $\mathfrak K$，则根据该定理，它必为共轭分量之和，因而是 $\mathfrak B$ 中的元素。

\paragraph{4. $\mathfrak K$ 的互补基及其同单位元分量的关系。}
从不变结构转到基表示时，这一结构把各个基归并为互补基对：

\paragraph{定理。}
\emph{$\mathfrak K$ 的每个 $P$-基 $t_1,\ldots,t_n$，都对应着 $\mathfrak K$ 中的一个互补基 $T_1,\ldots,T_n$，而后者的互补基又是原来的 $t$ 基。$K^{(i)}$ 中同互补基 $T_1,\ldots,T_n$ 同构的基 $A_1^{(i)},\ldots,A_n^{(i)}$，被定义为单位元分量 $e^{(i)}$ 在由 $t$ 给出的基表示
\[
e^{(i)}=A_1^{(i)}t_1+\cdots+A_n^{(i)}t_n
\]
中的系数。若 $\alpha_1^{(i)},\ldots,\alpha_n^{(i)}$ 是与各 $t_i$ 对应的基，则 $\alpha$ 与 $A$ 的矩阵互逆。若 $\mathfrak K$ 的两个基 $(s)$ 与 $(t)$ 通过变换
\[
(s)=(t)P
\]
相联系，则相应互补基之间有逆变关系
\[
(S)=P^{-1}(T).
\]}

\paragraph{证明。}
设 $e^{(i)}=A_1^{(i)}t_1+\cdots+A_n^{(i)}t_n$ 是单位元的 $n$ 个分量 $e^{(i)}$ 的基表示，其中取定了 $\mathfrak K$ 的一个 $P$-基 $t_1,\ldots,t_n$；该基于是也成为 $\mathfrak K_\Gamma$ 的 $\Gamma$-基。根据第 2、3 小节，$A_1^{(i)},\ldots,A_n^{(i)}$ 属于 $K^{(i)}$，而不同的 $e^{(i)}$ 对应于彼此共轭的系数系 $A$。若 $\alpha_1^{(i)},\ldots,\alpha_n^{(i)}$ 表示 $K^{(i)}$ 中与 $t_1,\ldots,t_n$ 同构对应的基，则还有以下关系：
\begin{equation}
e^{(i)}[t\to\alpha^{(i)}]=e;
\qquad e^{(i)}[t\to\alpha^{(j)}]=0\quad\text{当 }i\ne j. \tag{1}
\end{equation}
因为 $\mathfrak B_{K^{(i)}}^{(i)}$ 包含元素 $t_1-\alpha_1^{(i)},\ldots,t_n-\alpha_n^{(i)}$——它们构成 $\mathfrak B_{K^{(i)}}^{(i)}$ 的一个非独立 $P$-模基——所以在代入 $t=\alpha^{(i)}$ 时该理想消失。于是由
\[
e\equiv e^{(i)}\pmod{\mathfrak B_{K^{(i)}}^{(i)}}
\]
可得 (1) 的第一个关系；第二个关系则由 $i\ne j$ 时 $e^{(i)}$ 属于所有 $\mathfrak B_{K^{(j)}}^{(j)}$ 而得到。按系数展开后，关系式 (1) 可写成矩阵形式：
\begin{equation}
\begin{pmatrix}
\alpha_1^{(1)}&\cdots&\alpha_n^{(1)}\\
\vdots&&\vdots\\
\alpha_1^{(n)}&\cdots&\alpha_n^{(n)}
\end{pmatrix}
\begin{pmatrix}
A_1^{(1)}&\cdots&A_1^{(n)}\\
\vdots&&\vdots\\
A_n^{(1)}&\cdots&A_n^{(n)}
\end{pmatrix}=E_n. \tag{2}
\end{equation}
其中 $E_n$ 表示单位矩阵；由此，$A$ 已由 $\alpha$ 唯一确定。关系式 (2) 表明两个矩阵的行列式都不为零。它还表明，$A_1^{(i)},\ldots,A_n^{(i)}$ 关于 $P$ 线性无关；因为若存在一个线性相关关系，它也会对各共轭系成立，从而降低 $A$ 矩阵的秩，这与 $E_n$ 的秩为 $n$ 相矛盾。因此，每组 $A_1^{(i)},\ldots,A_n^{(i)}$ 都构成 $K^{(i)}$ 的一个基；这些基彼此共轭，所以它们都同 $\mathfrak K$ 中的同一个基 $T_1,\ldots,T_n$ 同构，而根据第 3 小节，$T_\lambda$ 由
\[
T_\lambda=A_\lambda^{(1)}e^{(1)}+\cdots+A_\lambda^{(n)}e^{(n)}
\]
确定。称基 $T_1,\ldots,T_n$ 为 $t_1,\ldots,t_n$ 的互补基，相应地，称 $A_1^{(i)},\ldots,A_n^{(i)}$ 为 $\alpha_1^{(i)},\ldots,\alpha_n^{(i)}$ 的互补基。由于 $A$ 已通过 (2) 唯一确定，而这一关系——通过转到转置矩阵——又从 $T$ 的基、亦即从 $A$ 确定 $\alpha$，所以原来的 $t$ 基就是 $T$ 基的互补基。因此，单位元各分量所给出的对应把所有基都分成互补基对。若 $s,S$ 是另一对互补基，满足
\[
e^{(i)}=B_1^{(i)}s_1+\cdots+B_n^{(i)}s_n
\qquad\text{且}\qquad (s)=(t)P,
\]
则由
\[
A_1t_1+\cdots+A_nt_n=B_1s_1+\cdots+B_ns_n
\]
——在代入 $(s)=(t)P$ 后作为关于 $t$ 的恒等式——可知，$B,A$ 的变换与 $s,t$ 的变换互为逆变。因而由同构可知，$S$ 和 $T$ 也有同样的关系。至此，该定理的各部分全都得证。

\paragraph{注。}
若 $c=C_1t_1+\cdots+C_nt_n$ 是 $\mathfrak K$ 中一个元素关于某个 $P$-基的表示，而 $c=c_1T_1+\cdots+c_nT_n$ 是同一元素关于互补基的表示，则
\[
C_\lambda=\operatorname{Sp}(cT_\lambda),
\qquad c_\lambda=\operatorname{Sp}(ct_\lambda),
\]
其中 $\operatorname{Sp}(x)$ 照例表示各共轭之和：
\[
\operatorname{Sp}(x)=\xi^{(1)}+\xi^{(2)}+\cdots+\xi^{(n)}
\]
对 $\mathfrak K$ 中任意 $x$ 都如此。事实上，
\[
c=\gamma^{(1)}e^{(1)}+\cdots+\gamma^{(n)}e^{(n)}
=\sum_\lambda\gamma^{(1)}A_\lambda^{(1)}t_\lambda+
\cdots+\sum_\lambda\gamma^{(n)}A_\lambda^{(n)}t_\lambda,
\]
由此比较系数便得到 $C_\lambda=\operatorname{Sp}(cT_\lambda)$。第二个关系可同样通过交换 $t$ 与 $T$ 得到。\srcfn{19)}{本小节的讨论表明，不变结构定理实质上比非不变的基定理简单。这里只是为了说明它们同已有结果的联系，才使用后者。事实上，伽罗瓦扩张环的分解以及单位元分量的引入，已经完成了互补基（下一节中的互补模）所完成的一切。下一节中，不同与互补模之间的联系也只是为了归入已有理论才给出。对第 5 小节中拉格朗日插值公式的讨论同样如此。}

\paragraph{5. 以定义方程为基础。拉格朗日插值公式。}
[编者注：本小节在手稿中仍未填写。]
% Source: Paper 43, printed pp. 17--21 (journal scan pp. 706--710)
